% Standalone article. Compile twice with pdfLaTeX; no external bibliography or figures.
\documentclass[a4paper,11pt]{article}
\usepackage[T1]{fontenc}
\usepackage[margin=23mm,headheight=15pt]{geometry}
\usepackage{amsmath,amssymb,amsthm,mathtools,bm,mathrsfs}
\usepackage{booktabs,longtable,array,enumitem,xurl}
\usepackage[hidelinks,bookmarksnumbered]{hyperref}
\usepackage{fancyhdr}
\pagestyle{fancy}\fancyhf{}
\fancyhead[L]{Pointwise near-Einstein constructions}\fancyfoot[C]{\thepage}
\setlength{\parskip}{4pt}\setlength{\emergencystretch}{3em}
\allowdisplaybreaks[3]\setcounter{tocdepth}{2}\numberwithin{equation}{section}
\setlist{itemsep=2pt,topsep=4pt,parsep=0pt,leftmargin=2em}
\newtheorem{theorem}{Theorem}[section]
\newtheorem{proposition}[theorem]{Proposition}
\newtheorem{lemma}[theorem]{Lemma}
\newtheorem{corollary}[theorem]{Corollary}
\newtheorem{criterion}[theorem]{Conditional criterion}
\theoremstyle{definition}
\newtheorem{definition}[theorem]{Definition}
\newtheorem{problem}[theorem]{Open problem}
\newtheorem{remark}[theorem]{Remark}
\DeclareMathOperator{\Ric}{Ric}\DeclareMathOperator{\Rm}{Rm}
\DeclareMathOperator{\Vol}{Vol}\DeclareMathOperator{\tr}{tr}
\DeclareMathOperator{\Id}{Id}\DeclareMathOperator{\Div}{div}
\DeclareMathOperator{\Met}{Met}\DeclareMathOperator{\sech}{sech}
\DeclareMathOperator{\arsinh}{arsinh}\DeclareMathOperator{\TT}{TT}
\DeclareMathOperator{\Sym}{Sym}\DeclareMathOperator{\Out}{Out}
\DeclareMathOperator{\dist}{dist}
\newcommand{\R}{\mathbb R}\newcommand{\Z}{\mathbb Z}
\newcommand{\CP}{\mathbb{CP}}\newcommand{\op}{\mathrm{op}}
\newcommand{\dd}{\,d\mu_g}
\newcommand{\norm}[1]{\left\lVert#1\right\rVert}
\newcommand{\source}[2]{\par\smallskip{\small
  Proof source: \href{../references/rounds/#1}{\nolinkurl{#1}}, #2.}\par}
\newcommand{\base}[2]{\par\smallskip{\small
  Inherited source: \href{#1}{\nolinkurl{#1}}, #2.}\par}
\newcolumntype{L}[1]{>{\raggedright\arraybackslash}p{#1}}
\title{Pointwise near-Einstein constructions on fixed smooth manifolds\\
  \large Closed sheared topology, exact open models, and correction in all tensor modes}
\author{}
\date{21 September 2026}
\begin{document}
\maketitle
\begin{abstract}
We develop constructions measured by the global, unweighted Ricci
endomorphism defect. Two inherited mechanisms---marked $A_1$ resolution
using a corrected Eguchi--Hanson core and logarithmic Weyl interpolation,
and slow interpolation on a fixed surface bundle---are connected with
exact complete double-cylinder models, nonlinear repair obstructions, and
fixed closed gluings with integral shear and arbitrary mapping classes.
For the closed gluings we retain every Ricci block, prove the small-error
estimate, and record the volume and almost-parallel-curvature properties.
For the balanced surface bundle, a receiver-weighted inverse retaining the
whole closed fibre and its transfer modes gives a linear-size exclusion
neighborhood for arbitrary tensor perturbations. Complete two-ended
approximate metrics provide a separate test of exact connector obstructions.
We formulate the remaining all-mode problem through a Schur complement,
actual marked boundary response, and a Banach-space a posteriori criterion
whose hypotheses imply the full equation. These conditional tools do not
establish a new closed Einstein metric or exclude every Einstein metric
on the constructed manifolds. Positive and negative branches, actual smooth
topology, higher-dimensional extensions, and the distinction between exact
and approximate matching remain explicit throughout.
\end{abstract}
\tableofcontents
\clearpage

\section{The construction problem and its quantifiers}\label{sec:problem}
\subsection{The interval question and its two sign branches}\label{sec:quantifiers}
Let $Q_n$ be the following working question. Do there exist real constants
$c_n<C_n$, depending only on $n$, such that every closed connected smooth
$n$-manifold carrying a metric
\[
             c_ng\le\Ric_g\le C_ng
\]
also carries some smooth positive-definite Einstein metric
$h$, with $\Ric_h=\lambda h$? The output has no prescribed sign,
proximity, complex structure, fibration, or symmetry. The working statement
is fixed independently of the remaining historical comparison between the
1982 question page and later renumberings.

For $\sigma\in\{+1,-1\}$, set
\begin{equation}\label{eq:defect}
 \delta_\sigma(g)=\norm{g^{-1}\Ric_g-\sigma\Id}_{L^\infty,\op},
 \qquad e_\sigma(M)=\inf_{g\in\Met(M)}\delta_\sigma(g).
\end{equation}
Let $P_n^\sigma$ assert that a positive threshold depending only on $n$
for $\delta_\sigma$ implies the same unrestricted Einstein conclusion.
Then
\begin{equation}\label{eq:branches}
                  Q_n\quad\Longleftrightarrow\quad P_n^+\lor P_n^-.
\end{equation}
Indeed, for $\widehat g=ag$ the covariant Ricci tensor is unchanged and
its endomorphism is divided by $a$. Every nondegenerate interval contains
a nonzero interior point $\mu$ and a same-sign neighborhood; scaling a
metric close to $\operatorname{sign}(\mu)$ by $1/|\mu|$ puts it in that
interval. Conversely a sufficiently small interval centered at $+1$ or
$-1$ is an admissible interval.

An affirmative result in every dimension needs at least one successful
sign in each dimension. A complete disproof needs both signs to fail in
one and the same dimension, possibly on different manifolds or on changing
topologies. A fixed $M$ with $e_\sigma(M)=0$ supplies a branch counterexample
only after every Einstein metric on that same smooth $M$ is excluded.
Higher-dimensional construction gaps obstruct an all-dimensional
affirmative proof; they are not extra requirements on a complete
fixed-dimensional disproof.

\subsection{Approximation, correction, and global existence}\label{sec:levels}
A family with small defect, a solution of a projected correction equation,
and an exact Einstein metric are three different mathematical outputs.
Integral rigidity, construction, and global classification are current
research emphases, not a restriction on possible methods. An existence
argument, a new invariant, or a different equivalent formulation is
relevant whenever it returns to \eqref{eq:branches}.

Theorems recalled below retain the hypotheses of the archived proofs.
Algebraic deductions and conditional analytic criteria are proved where
stated; proposed extensions are identified as open problems.
The new closed sheared examples make
the distinction particularly sharp: a qualitative global theorem forces
every hypothetical exact output to have no local Killing field, while
the constructed inputs are locally almost symmetric.

\subsection{Norms, scale, and compactness}\label{sec:norms}
In dimension four write $s=\tr_g\Ric$ and
$E^\circ=\Ric-\frac{s}{4}g$. The norms of $W^\pm$ are the
Hilbert--Schmidt operator norms on $\Lambda^2_\pm$; $\Rm_{\rm full}$
uses the four-index tensor norm. On a closed oriented manifold,
\begin{align}
8\pi^2\chi&=\int_M\left(|W^+|^2+|W^-|^2+
          \frac{s^2}{24}-\frac{|E^\circ|^2}{2}\right)\dd,\label{eq:CGB}\\
12\pi^2\tau&=\int_M(|W^+|^2-|W^-|^2)\dd,\\
\int_M|\Rm_{\rm full}|^2\dd
 &=32\pi^2\chi+4\int_M|E^\circ|^2\dd.\label{eq:energy}
\end{align}
The Eguchi--Hanson curvature energy is $12\pi^2$ in curvature-operator
normalization and $48\pi^2$ in four-index normalization. For a metric
with $\Ric\approx-3g$,
\[
 \delta_-(3g)=\tfrac13\norm{g^{-1}\Ric_g+3\Id}_\infty.
\]
Below an operator norm of a covariant Ricci defect always means that
one index has first been raised with the actual metric.
Finite-$L^p$ smallness, a weighted norm, or small support volume cannot
replace \eqref{eq:defect}. A complete noncompact Einstein metric is
not a closed-manifold counterexample.

\subsection{What has changed in the construction problem}\label{sec:progress}
The baseline already supplies $A_1$ logarithmic interpolation and the
fixed Dehn-twist bundle $E_D$. The later work integrates an entire
nonlinear open-model source, enlarges an arbitrary-tensor closed repair
exclusion neighborhood, and produces fixed closed topologies with
nonzero integral shear. Each result must identify its smooth object,
global error norm, and precise information about exact Einstein output.
The transition between these results explains the current direction.
Exact local models make boundary flux and external response unavoidable.
The closed receiver then selects a weighted norm in which every relevant
transfer mode can be retained, even though maximum-norm modulation fails.
Two-ended approximations show how the defect infimum can vanish inside
a class with no exact symmetric connector. Returning to compact geometry,
the fixed sheared construction separates slow mapping-class variation
from a rapid shear transition at a collapsed throat. Its remaining exact
equation must combine all these lessons: complete modes, actual boundary
response, and a scale-aware nonlinear estimate.
For the latest objects, preserving circle symmetry in a closed exact
ansatz would remove every possible solution. Small nonsymmetric modes,
including exponentially weak transmission through long collars, therefore
belong to the correction problem.
\source{SYNTHESIS.md}{construction and global-classification interfaces}

\section{Marked \texorpdfstring{$A_1$}{A1} resolution and logarithmic Weyl matching}
\label{sec:A1}
\subsection{A fixed smooth resolution}\label{sec:A1-theorem}
\begin{theorem}[Inherited marked $A_1$ approximation]\label{thm:A1}
Let $(M_0,g_0)$ be a fixed closed oriented Einstein four-orbifold,
$\Ric_{g_0}=\sigma g_0$, with finitely many $A_1$ points. Fix equivariant
normal coordinates, end orientations, boundary markings, and the actual
smooth Eguchi--Hanson fillings. Suppose the smooth core correction satisfies
\[
 D\Ric_{g_{\rm EH}}\widetilde h=g_{\rm EH},\qquad
 \widetilde h=H_*+O_2(\rho^{-2}),
\]
and its quadratic asymptotic term and the exterior jet are in a common
radial gauge. Then the fixed resolution $M$ carries smooth positive
metrics $g_a$ with
\[
             \delta_\sigma(g_a)\le C/|\log a|\longrightarrow0.
\]
They converge in Gromov--Hausdorff topology to $(M_0,g_0)$, and their
rescaled cores converge to the prescribed Eguchi--Hanson metrics.
Constants depend on the fixed finite marked data.
\end{theorem}
The explicit core below supplies the stated hypothesis in the inherited
$A_1$ setting. This is not an assertion for arbitrary ALE models.
A local obstruction to exact desingularization leaves open distant
Einstein metrics, and the theorem alone does not identify the smooth
type associated with an arbitrary boundary marking.
\base{20260916_construction.tex}{the marked $A_1$ theorem and three-region proof}

\subsection{The nonzero Ricci source in the core}\label{sec:A1-core}
A bare scaled Ricci-flat bubble cannot work:
\[
 (a^2g_{\rm EH})^{-1}
   \bigl(\Ric_{a^2g_{\rm EH}}-\sigma a^2g_{\rm EH}\bigr)
       =-\sigma\Id .
\]
The inherited construction on the $\mathcal O(-2)$ disk bundle uses
the biaxial family
\begin{align}
c_\Lambda&=(1-\Lambda/4)^2-1,\\
G_\Lambda(r)&=\frac{\sqrt{1+c_\Lambda r^2}}{r^4}
 \int_1^r\frac{s^3(4-\Lambda s^2)}
 {(1+c_\Lambda s^2)^{3/2}}\,ds,\\
F_\Lambda&=(1+c_\Lambda r^2)G_\Lambda,\\
q_\Lambda&=F_\Lambda^{-1}dr^2+
 \frac{r^2}{4}(\sigma_1^2+\sigma_2^2+G_\Lambda\sigma_3^2).
\end{align}
Here the invariant one-forms and angle periods have the fixed bolt
normalization of the source; the expression is used only where the
square root and positive metric coefficients are defined. The equation
\[
r(1+c_\Lambda r^2)G_\Lambda'
 +(4+3c_\Lambda r^2)G_\Lambda=4-\Lambda r^2
\]
together with Bianchi gives $\Ric(q_\Lambda)=\Lambda q_\Lambda$.
Differentiating at $\Lambda=0$ and subtracting the gauge term generated
by $V=\frac5{72}(r^3-r^{-1})\partial_r$ gives $\widetilde h$.
Bolt smoothness is checked in genuine smooth coordinates, not by
examining a single singular polar coefficient.

If $k=g_{\rm EH}$ and $r=a\rho$, the corrected core
$G_a=a^2(k+a^2\sigma\widetilde h)$ satisfies
\begin{equation}\label{eq:core-error}
 \norm{G_a^{-1}\Ric(G_a)-\sigma\Id}_{\op}
                   \le C(a^2+r^2).
\end{equation}
An $O(a^2)$ relative metric correction produces an order-one original-scale
Ricci correction because Ricci contains two derivatives and one index is
raised at the shrinking scale.

\subsection{The full Weyl jet and the logarithmic annulus}\label{sec:A1-log}
Equal Ricci contractions of the inner and outer quadratic jets imply,
in the common radial gauge, that their difference is
\[
 K_{ij}=-\tfrac13C_{ikjl}x^kx^l,\qquad
 \tr K=\Div K=\Delta K=0,\qquad K(x,\cdot)=0 .
\]
The algebraic Weyl tensor $C$ ranges over the full ten-dimensional
space $\Sym^2_0(\Lambda^+)\oplus\Sym^2_0(\Lambda^-)$, not just
the biaxial subspace. For $t=\log r$,
\begin{equation}\label{eq:log-Weyl}
 D\Ric_e[\eta(\log r)K]
             =-\frac{\eta_{tt}+6\eta_t}{2r^2}K.
\end{equation}
To see the coefficient $6$, use the homogeneity-two identity
$x\cdot\nabla K=2K$ in the radial product rule for $\Delta(\eta K)$
in dimension four. The trace and divergence identities eliminate the
other linearized Ricci terms.

Choose $a\ll b=a^{1/2}\ll R=a^{1/4}\ll1$ and
$L=\log(R/(8b))$. A fixed profile rescaled to this logarithmic interval
has $\eta_t=O(L^{-1})$, $\eta_{tt}=O(L^{-2})$.
Since $K=O(r^2)$, \eqref{eq:log-Weyl} costs $O(L^{-1})$.
The separate core bound \eqref{eq:core-error}, ALE-tail matching,
exterior normal-coordinate remainder, and nonlinear annular estimates
complete the proof. The linear identity alone does not prove the theorem.

The first ALE obstruction can remain order one after raising an index
for a specified fixed matching. Logarithmic matching changes the matching
mechanism. These statements are compatible. Logarithmic cutoffs and a
pointwise Ricci-approximation problem already occur in Ozuch's work
\cite{OzuchII}; comparison must concern the full norm, marking, and
geometric limit.

\subsection{Extensions and retained restrictions}\label{sec:A1-extensions}
Collinear ADE periods require the full $3b_2$-dimensional obstruction
space. Type-T and Wahl descent additionally require covers, periods,
finite averaging, and kernel conditions; $b_2=0$ of a quotient is not
a proof that its linearized kernel vanishes. Constant-scalar-curvature
correction requires a weight separated from the actual endpoint near-kernel:
an estimate for $0<\mu<2$ does not extend automatically to $\mu=2$.
Quartic matching requires a new extension theorem and tail estimates.

Spherical suspensions, the chirality at both poles, and four-node product
quotients need actual equivariant or Kirby data to determine smooth
type \cite{OR}. The archived same-chirality applications and
negative-branch GH separation are retained as such; the earlier
$S^2\times S^2$ application is not counted again.
Constants $1/3$ for a specified logarithmic profile, $1/8$ in another
restricted optimization, and $1/6$ for a constant-scalar model belong to
different gauges, cutoffs, and scales. None is a universal sharp constant.
\base{20260912.tex}{Chapters 1--4: flux, ADE/type-T/Wahl, scalar correction, and smooth models}

\section{A fixed surface bundle: slow interpolation and cubic repair}
\label{sec:ED}
\subsection{Keeping the monodromy fixed}\label{sec:ED-topology}
Fix $F=\Sigma_p$, $B=\Sigma_q$, $p,q\ge2$, and a nonseparating
Dehn twist $D$. On a fixed fibre collar put
\[
 h_z=du^2+a(u)^2(ds+z\zeta'(u)\,du)^2,\qquad
 a(u)=a_0\cosh u,\quad 0\le z\le1.
\]
Every $h_z$ is hyperbolic with the same area form, and
$h_1=D^*h_0$. A nonintegral local coordinate shear need not extend to
an isotopy of the whole fibre; it cannot remove the monodromy.
On a base collar with $b(t)=\ell\cosh t$ and half-length $L\to\infty$,
set $z(t)=\eta(t/L)$ and
\[
          G_\ell=dt^2+b(t)^2d\theta^2+h_{z(t)}.
\]
The metric is product outside the collar and closes using the fixed $D$.
All parameters describe metrics on one smooth bundle $E_D$.

\subsection{Every Ricci component}\label{sec:ED-Ricci}
Write $j=\tanh u$, $k=\tanh t$, $Q=a\zeta'$, and $v=Qz'$.
In the orthonormal coframe
$(dt,b\,d\theta,du,a(ds+z\zeta'\,du))$ the exact defect is
\begin{equation}\label{eq:ED-matrix}
(\Ric+G_\ell)_{\rm ON}=
\begin{pmatrix}
-v^2/2&0&0&(v_u+2jv)/2\\
0&0&0&0\\
0&0&-v^2/2&-(v_t+kv)/2\\
(v_u+2jv)/2&0&-(v_t+kv)/2&v^2/2
\end{pmatrix}.
\end{equation}
Differentiating frame coefficients alone would miss the quadratic
diagonal terms in the covariant derivative of the shape operator.
The mixed terms are essential to the estimate.
Since $z'=O(L^{-1})$, $z''=O(L^{-2})$, the global unweighted error is
$O(L^{-1})$. Moreover,
\[
\Vol(G_\ell)=16\pi^2(p-1)(q-1),\qquad
\norm{\nabla^m\Rm}_\infty\le C_m,\qquad
\norm{\nabla^m(\Ric+G_\ell)}_\infty\le C_m/L .
\]
The injectivity radius collapses and the diameter grows. In the shear
region $K(e_0,e_3)=v^2/4>0$ wherever $v\ne0$, so strict-negative-
sectional-curvature correction theorems do not apply.
\base{20260912.tex}{Chapter 6, BUN:slow, BUN:ricci-matrix, and BUN:slow-geometric}

\subsection{Balanced collars and the meaning of cubic order}\label{sec:ED-balanced}
Pinch the fibre waist $\varepsilon$ and the base waist
$\beta=\varepsilon^2$. Balancing the mixed sources gives
\begin{equation}\label{eq:balanced}
 \delta_-(G_\varepsilon)=\frac{2}{\pi^4}\varepsilon^2+O(\varepsilon^4).
\end{equation}
Three scalar correction equations, with their actual tail estimates,
produce an $O(\varepsilon^2)$ tensor correction on the same closed
manifold and a defect $O(\varepsilon^3)$. It need not preserve the
original projection as a Riemannian submersion.
Global TT pairings and the conformal Ricci equation give bounds in
specified neighborhoods or conformal classes. Cubic optimality there
is not an all-metric lower bound for $e_-(E_D)$.
The baseline free-quotient construction also permits independent and
noncommuting mapping classes. Later multitwist examples refine its
group and representation information, not the underlying interpolation
mechanism.
\base{20260916_global.tex}{the bundle-input and charged-center sections}

\section{Exact complete double-cylinder metrics}\label{sec:exact-open}
\subsection{All bounded invariant directions and an actual nonlinear family}
\label{sec:exact-family}
On $X=(\R\times S^1)^2$, with fixed angle periods, let
\begin{gather*}
 g_0=dt^2+\cosh^2t\,dy^2+du^2+\cosh^2u\,dx^2,\\
 f=\sech^2t,\qquad r=\sech^2u,\qquad J(v)=\arctan(\sinh v),\qquad
 W=\cosh t\cosh u .
\end{gather*}
For the orthonormal coframe $(e^0,e^1,e^2,e^3)$, put
$T_1=(e^0)^2-(e^1)^2$, $S_1=e^0e^1+e^1e^0$ and similarly
$T_2,S_2$. Let $P_t=\partial_t^2+\tanh t\,\partial_t$.
The Bianchi-gauged operator $2L=\nabla^*\nabla-2\mathring R$
has all ten invariant tensor blocks:
\[
\begin{array}{c|c}
\text{two factor traces}&-P_t-P_u+2\\
T_1,S_1&-P_t-P_u+2-4f\\
T_2,S_2&-P_t-P_u+2-4r\\
\text{four mixed polarizations}&-P_t-P_u+2-f-r .
\end{array}
\]
The classification in the source allows arbitrary joint $(t,u)$
dependence and yields exactly
\begin{align}\label{eq:eight-kernel}
H[z]={}&f[(a_0+a_1J(u))T_1+(b_0+b_1J(u))S_1]\\
 &+r[(A_0+A_1J(t))T_2+(B_0+B_1J(t))S_2].\nonumber
\end{align}
These eight bounded fields are not $L^2$ on the full product.
In the invariant $L^2(W\,dt\,du)$ domain the source proves instead
$\langle Lv,v\rangle\ge\frac18\norm v_2^2$.

Define the exact hyperbolic cylinder
\[
h_{\alpha,\gamma}(v,\theta)
=\frac{\cosh^2v}{\sinh^2v+e^{-\alpha}}\,dv^2
 +(\sinh^2v+e^{-\alpha})(d\theta+\gamma\sech^3v\,dv)^2 .
\]
Let its first pair of parameters depend affinely on $J(u)$ and its
second pair affinely on $J(t)$ as in \eqref{eq:eight-kernel}.
Their sum is a reference $g[z]$ of determinant one in the background
frame. Define $H^m_\delta$ by $W^\delta k\in H^m(g_0)$ within
the invariant tensor subspace.

\begin{theorem}[Archived complete nonlinear family]\label{thm:open-family}
For sufficiently small $z\in\R^8$ there is an analytic correction
$k[z]\in H^6_{1/4}$, with higher regularity, such that
\[
 G[z]=g[z]+k[z],\qquad \Ric(G[z])=-G[z],\qquad
 \beta_{g[z]}k[z]=0,\qquad \norm{k[z]}_{H^6_{1/4}}\le C|z|^2.
\]
The metrics are smooth, complete, and $T^2$ invariant, and for each fixed $j$,
$|\nabla_0^jk[z]|\le C_j|z|^2W^{-3/4}$.
Uniqueness is in the stated correction space and gauge.
\end{theorem}
\begin{proof}[Analytic mechanism]
Set $\xi=J(t)$, $\eta=J(u)$, $p=\cos\xi$, $q=\cos\eta$.
The exact nonlinear residual has smooth frame coefficients on the
closed square and vanishes on all four faces. Hadamard division gives
$E[z]=pqR[z]$ with $R[z]=O(|z|^2)$ smooth. Thus the whole source,
not just its quadratic truncation, lies in $H^m_\delta$ for $\delta<1/2$.
Conjugation by $W^\delta$ gives
\[
\operatorname{Re}\langle L_\delta v,v\rangle
\ge(1/8-\delta^2)\norm v_2^2,
\]
using $|\nabla\log W|^2\le2$ and the combined skew-adjoint drift terms.
At $\delta=1/4$ this yields the inverse by coercivity and regularity.
Contraction solves
\[
\Ric(g[z]+k)+g[z]+k+
\delta^*_{g[z]+k}\beta_{g[z]}k=0
\]
in a ball of radius $C|z|^2$.
The Bianchi one-form $\omega$ is $L^2$ and solves
$(\nabla^*\nabla-\Ric)\omega=0$. Negative Ricci follows from
$C^2$ proximity; cutoff integration forces $\omega=0$.
\end{proof}
A moving reference is necessary: a constant cylinder length change
has a fixed-reference Bianchi error independent of the other radial
variable and therefore outside $L^2$, regardless of its quadratic size.
The eight parameters are marked boundary/kernel parameters, not eight
proved inequivalent moduli.
\source{round2_construction.md}{all blocks, bounded kernel, physical scales, and invariant gap}
\source{round3_construction.md}{nonlinear target space, inverse, moving gauge, contraction, and Bianchi removal}

\subsection{Charge, endpoint flux, and unchanged-core insertion}\label{sec:charge}
Write a genuine invariant Einstein metric as
\[
G=h+H_{ab}(d\theta^a+\mathcal A^a)(d\theta^b+\mathcal A^b),\quad
v=\sqrt{\det H},\quad F=d\mathcal A,\quad \mathfrak q=vH*_hF .
\]
With generators $(\partial_y,\partial_x)$ and quotient orientation
$dt\wedge du$, the mixed equation gives $d\mathfrak q=0$.
Neither $H$ nor $h$ is assumed diagonal. Put $c=\pi^2/2$. Then
\begin{equation}\label{eq:charge}
\int_{\R^2}F=c(-b_1,B_1)^{\mathsf T}
 =\mathsf M[z]\mathfrak q[z],\qquad
\mathsf M[z]=\int_{\R^2}v^{-1}H^{-1}dA_h>0,\qquad \mathsf M[0]=c\Id.
\end{equation}
The density is bounded by a constant times
$\operatorname{diag}(\sech^3t\sech u,\sech t\sech^3u)$.
The correction to $\mathcal A$ has boundary integral tending to zero,
also after fixed parameter derivatives, so Stokes uses actual end markings.
The finite twist energy is
\[
\mathcal E_{\rm tw}=\mathfrak q^{\mathsf T}\mathsf M\mathfrak q
=c\,\mathfrak q\cdot(-b_1,B_1)
=c^2(-b_1,B_1)\mathsf M^{-1}(-b_1,B_1)^{\mathsf T}.
\]
It vanishes exactly when $b_1=B_1=0$. Constant twists $b_0,B_0$
are removed by angular shifts built from $\int_0^v\sech^3w\,dw$,
but these alter the end markings. Real lift differences do not vanish
merely because their torus values agree.

At second order the two net receiver sources are
\[
I_1=\frac{\pi^2}{8}(a_1^2-3b_1^2-A_1^2-B_1^2),\qquad
I_2=\frac{\pi^2}{8}(A_1^2-3B_1^2-a_1^2-b_1^2).
\]
Homogeneous end shifts move opposite face fluxes without changing their
sum $-I_i$. \emph{If} the actual exterior response vanished, these
relations would force $b_1=B_1=0$ and $a_1^2=A_1^2$.
That premise is unproved on $E_D$. The exterior source, adjoint
detector defect, gauge terms, and physical periods must all be retained.
The matrix in \eqref{eq:charge} is a two-channel twist response, not
a full tensor Dirichlet-to-Neumann map.
\source{round3_matching_next.md}{conditional matching and the actual exterior detector}
\source{round4_construction.md}{nonlinear charge, finite scattering matrix, and twist energy}

For $(b_1,B_1)\ne0$, the archived exact-core theorem excludes any
unchanged nonempty open piece of $G[z]$ in a connected closed Einstein
four-manifold. Escaping fixed-radius balls on the complete cover approach
product geometry with $\nabla\Rm\to0$. A cocompact isometry group would
bring them into a compact region, force $\nabla\Rm=0$ on an open set,
and hence everywhere by analyticity. The resulting
$\mathbb H^2\times\mathbb H^2$ geometry has integrable orthogonal
complement to every locally free two-dimensional abelian Killing action,
contradicting nonzero charge. Analytic continuation of a local isometry
between complete simply connected analytic covers transfers the closed
insertion to this contradiction. This argument neither gives a
quantitative distance from exact metrics nor excludes changes on every
open set.
\source{round4_exact_core.md}{analytic continuation and cocompactness}

\section{An arbitrary-tensor repair obstruction on the closed bundle}
\label{sec:closed-repair}
\subsection{The linear-size exclusion ball}\label{sec:exclusion}
\begin{theorem}[Archived closed repair theorem]\label{thm:exclusion}
For the balanced $G_\varepsilon$ on the fixed $E_D$, with the specified
single nonseparating fibre degeneration, there are $c_*,c_1,\varepsilon_0>0$
such that every smooth symmetric tensor $h$ satisfying
$\norm h_{C^2(G_\varepsilon)}\le c_*\varepsilon$ and
$0<\varepsilon<\varepsilon_0$ gives a positive metric with
\[
                \delta_-(G_\varepsilon+h)\ge c_1\varepsilon^3.
\]
There is no symmetry, support, block, or gauge restriction. If
$\norm h_{C^2}=o(\varepsilon)$, then
\[
\liminf_{\varepsilon\to0}\varepsilon^{-3}
\delta_-(G_\varepsilon+h)\ge\frac1{16\pi^4(p-1)}.
\]
\end{theorem}
Reducing the radius by a fixed factor and rescaling also excludes nearby
Einstein metrics with other Einstein constants. The inherited cubic
construction lies in this ball, so the order is sharp within the stated
neighborhood. The radius tends to zero; no distant Einstein metric is
excluded.

\subsection{The receiver measure, full kernel, and error ledger}\label{sec:weighted}
Let $E(g)=\Ric_g+g$. The actual global TT receiver $Q_\varepsilon$
obeys, with $\beta=\varepsilon^2$,
\begin{align}
\beta^{-1}\int\langle E(G_\varepsilon),Q_\varepsilon\rangle
 &=-c_0\varepsilon^3(1+o(1)),&c_0&=\frac1{2\pi^2},\\
\beta^{-1}\norm{J_GQ_\varepsilon}_1&\le C\varepsilon^2,&
\beta^{-1}\int|Q_\varepsilon|_*&\longrightarrow M_0=8\pi^2(p-1).
\end{align}
The nuclear norm $|\cdot|_*$ is dual to the operator norm.
Base-TT quadratic cancellation alone does not control the sign of
arbitrary quadratic responses. Passing to the whole fibre preserves the
thick region and both cusp limits, which a double-cylinder limit could lose.

After angular averaging, use the comparison product with the actual
closed fibre $F_\varepsilon$, per unit base-circle coordinate length, and
\[
d\nu=\sech t\,dt\,dA_{F_\varepsilon},\qquad
\norm U_{H^j_\nu}^2=\sum_{i=0}^j\int|\nabla^iU|^2d\nu .
\]
Its mass is $4\pi^2(p-1)$, uniformly in $\varepsilon$.
For smooth invariant $V\in H^2_\nu$, the archived estimate is
\begin{equation}\label{eq:weighted-inverse}
V=K+r,\qquad K=f(aT+bS)+A+J(t)B,\quad
A,B\in\TT(F_\varepsilon),\qquad
\norm r_{H^2_\nu}\le C\norm{LV}_{L^2_\nu}.
\end{equation}
All product geometric derivatives are included. The scalar gap
$\lambda_1(F_\varepsilon)\ge\lambda_*>0$ uses connectedness of the thick
complement for a single nonseparating pinch; it does not extend to an
arbitrary separating degeneration \cite{Ji}.

For completeness, with nonnegative fibre Laplacians and
$P=\partial_t^2+\tanh t\,\partial_t$, the blocks of $2L$ are
\[
\begin{array}{c|c}
\text{factor traces}&-P+\Delta_F+2\\
\text{horizontal trace-free}&-P+\Delta_F+2-4f\\
\text{mixed one-form coefficients}&-P+\Delta_{H,F}+2-f\\
\text{fibre trace-free}&-P+A_F .
\end{array}
\]
Here $A_F=2P_F\delta_F$ has nonzero spectrum at least two.
For a scalar coefficient,
\[
\int\sech t\, u(-P+V_0)u
=\int\sech t\,\bigl((u')^2+(V_0+2f-1)u^2\bigr).
\]
The resonant horizontal block factors through
\[
D_1=\partial_t+2\tanh t,\quad D_2=-\partial_t+\tanh t,\qquad
D_2D_1=-P+2-4f,\quad D_1D_2=-P+2-f .
\]
The partner is coercive; the kernel of $D_1$ is spanned by $f$.
The vertical TT equation $-Pu=F$ has homogeneous solutions $A+JB$,
and its zero-data integral inverse is
\[
r(t)=-\int_0^t\sech u\int_0^u\cosh s\,F(s)\,ds\,du .
\]
After conjugating by $\sqrt{\sech t}$ its kernel has integrable
exponential decay on each half-line. These estimates prove
\eqref{eq:weighted-inverse} on the full weighted domain, with exhaustion
justifying integrations at infinity. The transfer $JB$ is retained,
without claiming that it integrates to a global bundle parameter.

The kernel has a signed quadratic receiver:
\begin{equation}\label{eq:signed}
\mathcal B(K,K):=\frac12\int\langle D^2E[K,K],Q_0\rangle
             =-\frac{\pi}{8}\norm B_{L^2(F_\varepsilon)}^2\le0 .
\end{equation}
A determinant-preserving fibre path makes its horizontal Ricci
variation minus the squared fibre velocity. Together with
$|\mathcal B(U,V)|\le C\norm U_{H^2_\nu}\norm V_{H^2_\nu}$,
this bounds only the positive response needed to cancel the source.
A uniform maximum-norm inverse is neither asserted nor needed.
\source{round6_weighted_inverse.md}{weighted domains, partner operators, fibre TT transfer, gauge inverse, and receiver}

Set $a=\norm h_{C^2}$ and
$\eta=\delta_-(G_\varepsilon+h)+\varepsilon^2+\varepsilon a+a^2$.
The finite marked collar, cutoff, angular averaging, and product
comparison give forcing $O(\eta)$ and
$\norm{V_{\rm osc}}_{H^1_\nu}^2\le C\beta a^2$.
If $H=V+\mathcal L_Xg_*$ and $F=DE(H)$, naturality gives
\[
D^2E[H,H]-D^2E[V,V]=2\mathcal L_XF-DE[W].
\]
Only the second term disappears against the exact receiver. The first
costs $Ca\eta$. Including the exterior mass and every same-order term,
\begin{align}\label{eq:ledger}
\beta^{-1}\int\langle E(G+h),Q_\varepsilon\rangle
\le{}&-c_0\varepsilon^3(1+o(1))\\
 &+C[\varepsilon^2a+a\eta+\eta^2+
                 (\varepsilon+\beta)a^2+a^3].\nonumber
\end{align}
To prove Theorem~\ref{thm:exclusion}, first separate the case where the
desired lower bound already holds. Otherwise $\delta_-$ is at most
a fixed small multiple of $\varepsilon^3$, and
$a\le c_*\varepsilon$ makes the positive side of \eqref{eq:ledger}
smaller than half the source after choosing $c_*$ once.
Duality with $M_0$ yields the contradiction. For $a=o(\varepsilon)$,
all ledger terms are $o(\varepsilon^3)$ and the coefficient is $c_0/M_0$.
\source{round6_construction.md}{the finite-collar reduction and complete error ledger}
\source{round6_receiver_review.md}{the nonzero gauge forcing and scalar-gap hypotheses}

\subsection{What would enlarge the exclusion theorem}\label{sec:larger-repairs}
A global TT modulation counterexample has unit maximum size,
residual $\Theta(1/|\log\varepsilon|)$, and fixed positive distance
from reference tangents plus arbitrary infinitesimal diffeomorphisms.
It disproves an overstrong uniform maximum-norm inverse; its much larger
residual does not defeat Theorem~\ref{thm:exclusion}.
At $a=C\varepsilon$ for arbitrary fixed $C$, several cubic-order
terms in \eqref{eq:ledger} compete. Extending the result to every such
$C$, to an arbitrary $o(1)$ neighborhood, or to another background
requires new signed information. Conversely, changes beyond this
radius may be a route toward an exact metric. Neither outcome is
predetermined.
\source{round5_modulation.md}{the maximum-norm counterexample and its scale}

\section{Fixed closed gluings with mapping classes and integral shear}
\label{sec:sheared}
\subsection{The marked smooth object and its approximation theorem}
\label{sec:sheared-theorem}
Fix finitely many compact oriented surfaces $B_v$ with nonempty boundary
and $\chi(B_v)<0$, and closed oriented surfaces $F_v$ of a common genus
$p\ge2$. Pair every boundary in a connected finite graph, using fixed maps
\begin{equation}\label{eq:gluing}
\Psi_e(\phi,x)=(-\phi+u_e(x),f_e(x)),\qquad
u_e:F_v\to\R/(2\pi\Z),\quad f_e:F_v\to F_w.
\end{equation}
Each $f_e$ preserves orientation; its mapping class may have infinite
order and need not commute with any other one or extend over a block.
The resulting closed smooth oriented $M$ is fixed before $\ell$ varies.

\begin{theorem}[Archived fixed closed gluing]\label{thm:sheared}
The manifold $M$ carries smooth positive metrics $G_\ell$, $\ell\downarrow0$,
such that
\begin{align}
\norm{G_\ell^{-1}\Ric_{G_\ell}+3\Id}_\infty
 &\le C(|\log\ell|^{-1}+\ell),\label{eq:sheared-bound}\\
\Vol(G_\ell)&=\frac{4\pi^2}{9}\chi(M),\qquad
\norm{\nabla^j\Rm_{G_\ell}}_\infty\le C_j .
\end{align}
Every fixed covariant derivative of the defect has the same small
order bound. The marked gluing maps and integral shear classes do not
vary. For these families one also has
$\norm{\nabla^j\Rm_{G_\ell}}_\infty\to0$ for every $j\ge1$.
Consequently $e_-(M)=0$ after scaling by three, with volume
$4\pi^2\chi(M)$ in that normalization.
\end{theorem}

\subsection{Long collars absorb the mapping class}\label{sec:path-blocks}
Give each factor curvature $-3$ and every base boundary length $\ell$.
The hyperbolic collar theorem \cite{Rupflin} gives
\[
a_\ell(t)=\frac{\ell}{2\pi}\cosh(\sqrt3t),\qquad
L_\ell=\frac1{2\sqrt3}\arsinh
           \frac1{\sinh(\sqrt3\ell/2)}\asymp|\log\ell|.
\]
In a signed collar $[-L_\ell,L_\ell]$, the right product pulls back to
\[
dt^2+a_\ell^2(d\phi+\beta_e)^2+f_e^*h_w,\qquad
\beta_e=-u_e^*d\phi .
\]
The closed one-form $\beta_e$ has periods in $2\pi\Z$ and need not be exact.
Choose a fixed smooth path of hyperbolic metrics $h_z$ from $h_v$
to $f_e^*h_w$, constant near its endpoints, and traverse it on
$[-L_\ell,-1]$ with $z'=O(L_\ell^{-1})$. The metric there is
$g=dt^2+a_\ell^2d\phi^2+h_{z(t)}$. No global isotopy to $f_e$ is assumed.

Put $S=\frac12h^{-1}\dot h$, $H=\tr S$, $k=a_\ell'/a_\ell$,
and $U=a_\ell^{-1}\partial_\phi$. Gauss--Codazzi gives every block of
$E=\Ric+3g$:
\begin{align}
E_{tt}&=-\dot H-\tr S^2,& E_{UU}&=-kH,\label{eq:path-Ricci}\\
E_{FF}^{\sharp}&=-D_tS-(k+H)S,&
E_{tF}&=\Div_h S-dH,\qquad E_{tU}=E_{UF}=0.\nonumber
\end{align}
In particular a nonzero fibre-area trace cannot be dropped.
The fixed compact path gives
$S,H,\Div S,dH=O(L_\ell^{-1})$ and
$D_tS,\dot H=O(L_\ell^{-2})$, while $|k|\le\sqrt3$.
The full defect is $O(L_\ell^{-1})$. Coefficients are circle invariant;
the circle connection uses $k$, and there is no hidden $\ell^{-1}$ loss.
The path need not preserve its pointwise area form. Its total fibre
area is fixed by Gauss--Bonnet.

\subsection{The collapsed throat absorbs a nonexact integral shear}
\label{sec:shear-blocks}
After the fibre path has ended, fix $h=f_e^*h_w$ and on $[0,1]$ use
\[
g=dt^2+a_\ell^2(d\phi+z_s\beta_e)^2+h,
\]
where $z_s$ is a fixed smooth step, constant near both endpoints.
Set $b=\beta_e$, $D=\Div_hb$ (divergence of the dual vector field).
For horizontal fibre vectors $X,Y$, every nonzero error block is
\begin{align}
E_{tt}&=-a^2(z_s')^2|b|^2/2,&
E_{UU}&=a^2(z_s')^2|b|^2/2,\label{eq:shear-Ricci}\\
E_{XY}&=-a^2(z_s')^2b(X)b(Y)/2,&
E_{tU}&=az_s'D/2,\nonumber\\
E_{XU}&=-a(z_s''+3kz_s')b(X)/2,&E_{tX}&=0.\nonumber
\end{align}
The coefficient $3k$ is the derivative of $a^3$ in the mixed
connection equation, and is necessary even when $b$ is harmonic.
In the actual orthonormal splitting,
\[
\norm E_\op\le
\frac{a^2(z_s')^2|b|^2}{2}
\frac a2\sqrt{(z_s')^2D^2+(z_s''+3kz_s')^2|b|^2}.
\]
Since $a=O(\ell)$ on this fixed interval, the full error is $O(\ell)$.
Outside it the connection is flat. At the right endpoint the metric is
the actual pullback by \eqref{eq:gluing}; a nonexistent global real
primitive of $\beta_e$ is never introduced.

\subsection{Assembly, volume, almost parallel curvature, and both orientations}
\label{sec:assembly}
The two operations have disjoint supports. Their profiles are constant
near every interface, so all jets match. The finite number of disjoint
collars makes the global bound the maximum of the regional bounds,
proving \eqref{eq:sheared-bound}. The shear has determinant one;
the fibre paths have fixed total area. Consequently
\[
\Vol(G_\ell)=\sum_v\operatorname{Area}(B_v)\operatorname{Area}(F_v)
=\frac{4\pi^2}{9}\sum_v\chi(B_v)\chi(F_v)
=\frac{4\pi^2}{9}\chi(M).
\]
The gluing hypersurfaces have zero Euler characteristic.

Let $\mathcal R_0$ be the algebraic curvature tensor of two
curvature-$-3$ planes in the local actual splitting, and $P$ its first
plane projection. With $\epsilon_\ell=L_\ell^{-1}+\ell$,
\[
|\nabla^j(\Rm-\mathcal R_0)|\le C_j\epsilon_\ell\quad(j\ge0),\qquad
|\nabla^jP|\le C_j\epsilon_\ell\quad(j\ge1).
\]
The only mixing coefficients are $S=O(L_\ell^{-1})$ and
$az_s'\beta_e=O(\ell)$. Since $\mathcal R_0$ is algebraic in $g$ and $P$,
\[
-3-C\epsilon_\ell\le\sec\le C\epsilon_\ell,\qquad
\norm{\nabla^j\Rm}_\infty\le C_j\epsilon_\ell\quad(j\ge1).
\]
Collapse and growing diameter remain possible. Almost parallel curvature
does not imply the existence of a locally symmetric metric on $M$.

One can additionally make these metrics almost K\"ahler in both
orientations. Isotope $f_e$ to preserve chosen equal-total-area forms;
the boundary isotopy extends over a collar and preserves smooth type.
For a fibre path with endpoint area form $\omega$, choose
$d\sigma_z=\omega_z-\omega$ through a fixed linear primitive operator.
Then $\sigma_0=\sigma_1=0$. Moser's argument in an auxiliary parameter
on $(1-s)\omega+s\omega_z$ gives identity pullbacks at both endpoints,
avoiding spurious endpoint holonomy. The common-area path admits
\[
\Omega_\pm=a_\ell dt\wedge(d\phi+z_s\beta_e)\pm\omega,\qquad
\Omega_\pm^2=\pm2a_\ell dt\wedge(d\phi+z_s\beta_e)\wedge\omega\ne0.
\]
These closed forms match the product forms and are metric-compatible.
Complex integrability is not asserted.
\source{round11_closed_gluing.md}{complete Ricci blocks, fixed markings, volume, curvature, and endpoint-fixed Moser construction}

\subsection{Which gluings give new smooth objects}\label{sec:new-topology}
Identity mapping classes, and some finite-isometry holonomies, can return
to old surface bundles after interchanging the surface roles.
Instead, self-glue $F_{h,2}\times\Sigma_g$, $h\ge1$, $g\ge2$, with
a nonzero integral shear $u$ and an infinite-order mapping class $f$.
Writing $A=\pi_1(F_{h,2})$, $\Gamma=\pi_1(\Sigma_g)$, the actual group is
\[
\left\langle A\times\Gamma,t\ \middle|\
tc_1t^{-1}=c_2^{-1},\
t\gamma t^{-1}=c_2^{u(\gamma)}f_*(\gamma)\ (\gamma\in\Gamma)\right\rangle .
\]
The archived group theorem excludes every extension of one closed
hyperbolic surface group by another. Asphericity and positive Euler
characteristic eliminate the remaining oriented closed surface-bundle
possibilities. Thus the specified examples genuinely lie outside that
class. They have $\chi=4h(g-1)>0$ and $\tau=0$.
Their walls have nonamenable group $\Z\times\Gamma$ and the pieces are
products of hyperbolic surfaces; obstruction theorems for infranil-fibre
higher graph manifolds do not automatically apply.
\source{round11_graph_topology.md}{the HNN presentation, nonextension, and no-surface-bundle theorem}

\section{Complete approximate connectors and exact exclusions}\label{sec:connectors}
\subsection{Local charged centers and a one-ended charge release}\label{sec:one-ended}
The baseline already contains stable analytic Koiso germs with nonzero
charge, all Cauchy constraints, and genuine periodic torus data.
The later explicit parametrization uses
$M(\theta)=R_\theta\operatorname{diag}(e^a,e^{-a})R_\theta^{\mathsf T}$,
$Q_0=\mathfrak q^{\mathsf T}M^{-1}\mathfrak q/2$, and
$b=2\sinh a/\sqrt{12+2Q_0}$.
The $\Ric=-3g$ evolution has central $H=0$, $H'=3-Q_0$,
positive orbit-plane curvature $3+Q_0/2$, and a strictly stable
center when $\max Q_0<3$. This is analytic local existence, not a
smooth elliptic Cauchy well-posedness theorem or a complete continuation.

A different complete metric on $D^2\times T^2$ releases charge.
For a fixed cutoff $\chi$ equal to one on $[0,1]$ and zero on $[2,\infty)$,
put
\[
m_R(r)=\chi(r/R),\quad V_R=r^2-m_R/r,\quad
g_R=\frac{dr^2}{V_R}+V_Rd\vartheta^2+r^2(dx^2+dy^2),\quad r\ge1.
\]
Period $4\pi/3$ for $\vartheta$ makes the bolt smooth.
The inner region is an exact toral Einstein black hole and the far
region is exactly hyperbolic. The error eigenvalues are
$m_R''/(2r),m_R''/(2r),m_R'/r^2,m_R'/r^2$, so the full defect
is $O(R^{-3})$. For integral generators
$\partial_x+c\partial_\vartheta,\partial_y$, the charge is
$(c(3m_R-rm_R'),0)$ and changes across the small-error transition.
This one-ended topology is not a two-ended torus connector.
\source{round6_global.md}{charged germs, support-volume limitation, and one-ended bolt construction}

\subsection{An explicit two-ended Kottler approximation}\label{sec:Kottler}
On the fixed $\R_u\times S^1_\phi\times\Sigma$, with
$K_{h_\Sigma}=-1$ and $\phi$ of period $2\pi$, put $y=d\cosh u$ and
\begin{equation}\label{eq:Kottler}
g_d=\frac{1+y}{3+y}\,du^2+
\frac{y^2(y+3)}{3(y+1)}\,d\phi^2+\frac{(1+y)^2}{3}h_\Sigma.
\end{equation}
For $0<d<1/2$ this is smooth and complete with two AH ends of fixed
conformal metric $d\phi^2+h_\Sigma$. Its exact defect is
\begin{equation}\label{eq:Kottler-defect}
\norm{g_d^{-1}\Ric_{g_d}+3\Id}_\infty
                 =\frac{d(2d+3)}{(1+d)^2}<3d.
\end{equation}
Indeed, with $r=(1+y)/\sqrt3$,
$F=r^2-1+2/(3\sqrt3r)$ and $q=F(1-d^2/y^2)$, the metric is
$dr^2/q+Fd\phi^2+r^2h_\Sigma$ doubled at $y=d$.
The Ricci-error eigenvalues are
\[
-\frac{3d^2}{y(1+y)^2},\qquad
\frac{d^2(2y+3)}{y(1+y)^2}\quad\text{with multiplicity three}.
\]
The second dominates and decreases for $y\ge d$, proving the
continuous supremum. A fixed cutoff in $d^2/y^2$ produces exact
Kottler exteriors with $O(d^2)$ transition error and $O(d)$ central
error. The physical error-support length is $O(\log(1/d))$.
Curvature derivatives of fixed order remain bounded; the smallest
circle radius is comparable to $d$.

The optical zero mode
$w=(1+y)^{3/2}/(y\sqrt{y+3})$, normalized by its maximum, has total
endpoint overlap $2d\sqrt{d+3}/(1+d)^{3/2}\to0$.
Thus the product of error and support length can tend to zero without
contradicting the zero mode: its boundary overlap also tends to zero.
The parameter derivative further gives a non-gauge approximate
infinitesimal deformation, with unit-scale distance from gauge and
$O(d)$ linearized error; this alone is not a weighted or boundary-fixed
inverse statement.
\source{round10_kottler.md}{the eigenvalues, exact-end cutoff, optical localization, and approximate deformation}

The exact free-circle connector theorem excludes a two-ended smooth
conformally compact negative Einstein metric within its stated
free-circle class. It is compatible with \eqref{eq:Kottler-defect}.
The boundary is $S^1\times\Sigma$, not $T^3$, and the statement is
not a free-$T^2$ counterexample or a closed result. Boundary circle
symmetry need not extend to interior isometries: the required relative
fundamental-group hypothesis fails for the disconnected boundary here.
Completely nonsymmetric fillings are not thereby excluded.
\source{round9_connector.md}{the exact free-circle boundary-connectivity theorem and its scope}

\subsection{Fixed nontrivial circle bundles}\label{sec:cusp-doubling}
Let $P_m\to\Sigma$ have Chern number $m$. Suppose an exact invariant
$\Ric=-3g$ source has one PE end and one cusp, with relative covariant
two-jet error $O(e^{-\beta r})$ from
\[
G_0=\tfrac13dr^2+c_0^2e^{-2r}\theta^2+h,\quad
h=\tfrac13h_\Sigma,\quad d\theta=f\,dA_h,\quad
f\operatorname{Area}(h)=2\pi m .
\]
The Li--Liu construction supplies the relevant surface-cusp sources
\cite{LiLiu}; the source hypothesis includes their precise asymptotics.
Join two copies through
$g_C=dt^2+a(t)^2\theta^2+h$, $a=d\cosh(\sqrt3t)$.
Its error eigenvalues are
$0,a^2f^2/2,-a^2f^2/2,-a^2f^2/2$, with zero mixed components.
At a seam $a(T)=a_0$, the relative two-jet mismatch with the exponential
cusp is $O(d^2/a_0^2)$. Nonlinear connection-difference estimates give
\begin{equation}\label{eq:cusp-bound}
\delta_{-3}(g_{a_0,d})\le
 C(a_0^\beta+a_0^2+d^2/a_0^2),
\end{equation}
where $\delta_{-3}(g)=\norm{g^{-1}\Ric_g+3\Id}_\infty$.
Taking $d=a_0^2$ yields complete approximations on the same
$\R\times P_m$, with two unchanged PE ends.
The connection is global on the bundle; no global base potential is
required. Fixed Chern flux costs $a^2|F|^2$ when the circle collapses.
The trivial-bundle exact formula \eqref{eq:Kottler-defect} is not reused
for $m\ne0$. Flat $T^3$ AH cusps and ACH Nil cusps have different
curvature balances and are not covered by this substitution.
\source{round10_cusp_doubling.md}{source hypotheses, global connection, relative seam estimates, and cusp-type limits}
\source{round10_connector.md}{connector comparisons and charge leakage}

\section{Correction in all modes and actual matching}\label{sec:all-modes}
\subsection{The full projected equation and its Schur complement}\label{sec:Schur}
Fix the Einstein constant or include it among the unknowns, and fix a
geometric gauge. Write
$F_g(h)=r+Lh+N(h)$ between declared Banach spaces. For a closed smooth
background one possible local choice is the Bianchi-gauged metric
neighborhood in $C^{k+2,\alpha}(S^2T^*M)$ mapping to
$C^{k,\alpha}(S^2T^*M)$, $k\ge0$, $0<\alpha<1$.
On degenerating families a weighted or anisotropic norm may be preferable;
its domain, transmission conditions, and gauge variables must be specified.
For a boundary problem these conditions belong to the operator domain.

Let $X=X_0\oplus X_1$ and $Y=Y_0\oplus Y_1$ be topological direct
sums with bounded projections, and write $h=z+v$.
Both equations $\Pi_1F_g(z+v)=0$ and $\Pi_0F_g(z+v)=0$ are necessary.
The low-mode space must retain genuine kernels, degenerating near-kernels,
global parameters, and necessary frame and gauge variables. Dropping
the scalar kernel or relabeling TT transfer as gauge changes the problem.

\begin{criterion}[Banach-space Schur elimination]\label{thm:Schur}
Let $L:X\to Y$ be bounded with blocks $L_{ij}:X_j\to Y_i$.
Suppose $L_{11}:X_1\to Y_1$ is a bounded isomorphism with bounded inverse
in the stated spaces and boundary conditions. Then
$L(z+v)=-r$ is equivalent to
\begin{align}
(L_{00}-L_{01}L_{11}^{-1}L_{10})z
 &=-r_0+L_{01}L_{11}^{-1}r_1,\label{eq:Schur}\\
v&=-L_{11}^{-1}(r_1+L_{10}z).\nonumber
\end{align}
The whole $L$ is a bounded isomorphism if and only if the Schur
operator $X_0\to Y_0$ in \eqref{eq:Schur} is.
\end{criterion}
\begin{proof}
The second row uniquely determines $v$; substitution gives the first
row. Conversely the displayed pair solves both rows. The same
substitution constructs bounded inverses in either direction.
\end{proof}
This is an algebraic fact under an analytic isomorphism hypothesis,
not an established inverse theorem for the closed gluing. A nonlinear
reduction needs a $C^1$ solution $v(z)$ on a controlled domain and the
remaining obstruction $\mathcal O(z)=\Pi_0F_g(z+v(z))$.
At a solved complement equation,
$Dv=-(D_v\Pi_1F)^{-1}D_z\Pi_1F$, so $D\mathcal O$ is the
corresponding Schur complement. External response changes both the
effective source and the low-mode operator; neither $r_0$ nor
$L_{00}$ alone decides solvability.

\subsection{Angular coercivity and long-collar near-kernels}\label{sec:angular}
For a local circle radius $a(t)$, a nonzero Fourier mode $m$ has
principal cost $m^2/a^2$. This may improve inversion of high modes,
but says nothing by itself about the zero mode, long transfer modes,
connection couplings, negative potentials, or gauge. A complete method
needs estimates in mixed norms: control high modes through the principal
operator, retain low-mode propagation, and return exterior effects
through the actual boundary response.

For the self-glued manifolds of Section~\ref{sec:new-topology}, the
archived global result states that nonzero $u$ forces a discrete
universal-cover isometry group for every metric, and hence no local
Killing germ anywhere for any Einstein metric. Its proof uses the
specific group properties, universal-cover rigidity \cite{FW}, and
analytic continuation; it is not an arbitrary-tensor distance bound.
An exact solution must therefore have angular dependence in every
open region where the approximation has a circle symmetry, however
small that dependence may be.
\source{round12_symmetry.md}{complete local symmetry breaking, including finite covers}

\subsection{Both ends must be realized by one boundary problem}\label{sec:matching}
Actual matching includes the two induced metrics, second fundamental
forms, connection holonomy, constraints, periods, and the fixed
topological markings. The realizable Cauchy data form constrained
sets; two independent successful half-core constructions do not prove
their intersection. Ellipticity prevents arbitrary independent
prescription of all Dirichlet and Neumann data. A claim that several
endpoint values are adjustable needs an appropriately indexed boundary
problem, its kernel, and parameter derivative.

A useful first target is a finite-domain boundary response with a
proved principal symbol, cokernel, and parameter dependence, followed
by control as the truncation length grows. If its inverse grows
faster than the source decays, finite-domain solvability does not close
the gluing. The same discipline applies to graph blocks, multi-node
orbifolds, and noncompact cores, without forcing them into one model.

\subsection{An exact invariant identity and its escape directions}\label{sec:invariant-repair}
For an invariant collar
$g=dt^2+h_t+a^2\theta^2$, put $S=\frac12h^{-1}\dot h$,
$H=\tr S$, $d\theta=dt\wedge\beta+B$, and
$Q=|\beta|^2+\frac14|B|_{\rm tens}^2$.
The exact vertical-minus-normal Ricci identity is
\begin{equation}\label{eq:scalar-difference}
\Ric(U,U)-\Ric(T,T)
=H'-(a_t/a)H+|S|^2-\Delta_ha/a+a^2Q.
\end{equation}
Here $\Delta_h=\operatorname{div}\nabla$ in this formula. For any
Einstein constant the left side vanishes.
When $a=a(t)$ it gives the positive receiver
\[
\int_{-L}^L(|S|^2/a+aQ)\,dt
             =H(-L)/a(-L)-H(L)/a(L).
\]
Let $I_a=\int a\,dt$ and let $C_\partial$ be the maximum norm of
the endpoint expression. The marked Teichm\"uller displacement satisfies
$d_{\mathcal T}([h_-],[h_+])^2\le I_aC_\partial$.
For a radius comparable within factor $\Lambda$ to the cosh radius
and $|H(\pm L)|\le\eta$,
\[
\eta\ge\frac{\sqrt3}{4\Lambda^2}
                     d_{\mathcal T}([h_-],[h_+])^2 .
\]
The collapse scale cancels. Pseudo-Anosov translation length
$\log\lambda(f)>0$ remains an obstruction as the reference conformal
structure changes. This applies to an invariant unit-normal-foliation
sector, unlike Theorem~\ref{thm:exclusion}.

Spatial radius variation is a genuine additional term:
$\int_F(\Delta a)/a^2=2\int_F|\nabla a|^2/a^3$, not zero.
In the common-area sector $H\equiv0$, multiplying
\eqref{eq:scalar-difference} by $a$ and integrating over the closed
fibre forces $S=d\theta=0$; then $\Delta a=0$, so the metric is
locally a product of equal-curvature surfaces. Exact connection
transport contributes nonnegative receiver energy and cannot cancel it.
None of these identities identifies the cokernel of the unrestricted
Einstein problem.
\source{round12_repair.md}{exact connection transport, Teichm\"uller displacement, spatial radius, and common-area rigidity}

\section{Open construction directions}\label{sec:directions}
\subsection{Other ALE data and higher dimensions}\label{sec:other-ALE}
\begin{problem}
Choose an actual ADE, type-T, or noncollinear-period core; compute the
image of the quadratic source in every obstruction direction, prove
the solvability conditions and tails, and then identify the marked
global smooth filling. Determine separately whether unweighted Ricci
approximation, exact desingularization, and global Einstein existence hold.
\end{problem}
Yao's higher-dimensional obstruction \cite{Yao} supplies a comparison
under noncollapse, uniform energy, and the specified smooth Ricci-flat
ALE blow-up assumptions. It is not a classification of all
higher-dimensional resolutions. Huang--Ozuch's dimension-five gap
analysis \cite{HO} is another dimension-specific input.
The dimension and sign hypotheses must survive any transplantation.

\subsection{The positive branch and complete multi-node coupling}\label{sec:positive}
Marked resolutions of spherical and product quotients test $P_4^+$.
A finite-dimensional reduction must include sources, receivers,
same and opposite chirality, adjacency, diagonal terms, and self-action.
A proposed $t^3$ sign cannot be read through an $O(t^{5/2})$
remainder; prove $o(t^3)$ or include same-order logarithms explicitly.
Rigorous spectral tails and continuous-operator estimates take priority
over additional floating-point digits.
A genuine transverse zero would lead toward existence. A local nonzero
obstruction still leaves the unrestricted Einstein problem on the
smooth resolution. Mixed-chirality rows are not assumed unsolvable.
\base{20260916_construction.tex}{full-kernel and four-node coupling calculations}

\subsection{Simultaneous shear and mapping-class variation}\label{sec:network}
The present construction separates the two supports to close the error
estimate. Overlapping transitions, several noncommuting connection
directions, or a graph of responding nodes could enlarge the smooth
input class or alter the exact correction equation. The integral
shear and marking must remain fixed. If $f_\ell$ varies, determine
whether smooth type is constant or explicitly study a changing-topology
sequence.
\begin{problem}
For a fixed gluing graph, derive a nonlinear low-mode system with the
actual topological holonomy. Its edge variables must come from
realizable geometric boundary data. Decide whether this system is
only necessary or is equivalent to the full equation with controlled
complement and gauge.
\end{problem}

\subsection{Small symmetry breaking may carry decisive information}\label{sec:small-asymmetry}
The no-local-Killing theorem gives no lower bound on the size of
asymmetry. A completely nonsymmetric exact metric could in principle
use very small angular components to connect the local product regions.
Their transmission through a long collar may be exponential, while
their quadratic feedback may enter a receiver direction. Proving that
this cannot cancel the source would give quantitative rigidity;
constructing a cancellation would give a new exact mechanism.
Deleting a mode because it is visually small settles neither question.

\subsection{Changing the reference beyond the old neighborhood}\label{sec:beyond-ball}
On $E_D$ one may vary background moduli, conformal class, or a transfer
parameter, or examine metrics not close to $G_\varepsilon$.
For $h\sim C\varepsilon$ the small-$c_*$ theorem does not apply.
If a dangerous mode enters a new reference family, recompute the
global TT pairing, its measure, and its signs in that family; constants
from the old background cannot be reused without comparison.
Both an enlarged exclusion and a successful exact construction are
legitimate outcomes.

\section{A posteriori certification of the full equation}\label{sec:posteriori}
\begin{theorem}[Injective-preconditioner contraction]\label{thm:posteriori}
Let $X,Y$ be real Banach spaces, $U\subset X$ open,
$\overline B_X(x_0,r)\subset U$, and $F:U\to Y$ continuously
Fr\'echet differentiable. Let $A:Y\to X$ be bounded linear and
\emph{injective}. Suppose
\[
\norm{AF(x_0)}_X\le Y_0,\qquad
\norm{I_X-ADF(x_0)}_{\mathcal L(X)}\le Z_0,
\]
and for all $x,y$ in the closed ball,
\[
\norm{A(DF(x)-DF(y))}_{\mathcal L(X)}
                  \le Z_2\norm{x-y}_X .
\]
If
\begin{equation}\label{eq:radii}
Y_0+Z_0r+\tfrac12Z_2r^2\le r,\qquad Z_0+Z_2r<1,
\end{equation}
then the ball contains exactly one zero of $F$.
\end{theorem}
\begin{proof}
Set $T(x)=x-AF(x)$. Taylor's integral formula on the convex ball gives
$\norm{T(x)-x_0}\le Y_0+Z_0r+\frac12Z_2r^2\le r$.
Also
$\norm{DT(x)}\le Z_0+Z_2\norm{x-x_0}\le Z_0+Z_2r<1$.
Thus $T$ is a strict contraction of the complete closed ball into
itself. Its unique fixed point $x_*$ satisfies $AF(x_*)=0$.
Injectivity of $A$ gives $F(x_*)=0$. Every zero is a fixed point,
so uniqueness follows within the stated ball.
\end{proof}
The injectivity assumption cannot be dropped: contraction proves
$AF=0$, and a nonzero residual in $\ker A$ would remain invisible.
It suffices, for example, to construct $A$ as a bounded isomorphism;
alternatively a separate verified left bound on the relevant residual
space can establish injectivity. A numerical inverse of a finite
matrix proves neither property for the full function-space operator.

For an Einstein application one must additionally verify that the
certified ball consists of positive metrics, includes all continuous
spectral tails and boundary conditions, and has a valid gauge-removal
argument. A zero of a gauged equation is not automatically an Einstein
metric. On a closed negative-Ricci ball the usual Bianchi equation and
integration can remove the gauge; in other settings the appropriate
boundary or decay and kernel arguments are separate requirements.
Uniqueness is local to the chosen gauge and ball.

A practical candidate is a certified low-mode numerical inverse joined
to an analytic high-mode inverse, with rigorous bounds for
$Y_0,Z_0,Z_2$ and all coupling errors. Wang's construction
\cite{Wang} provides a specific external example where rigorous local
numerics, a fixed-point correction, and a separate complete AH
extension are all proved. It does not discharge the boundary, kernel,
or closed-topology conditions here. Failure of \eqref{eq:radii} may
reflect a poor candidate, a missing mode, a large inverse, or excessive
nonlinearity; it is not a nonexistence theorem.

\section{Two all-mode programs}\label{sec:programs}
\subsection{The complete low-energy system on a fixed sheared manifold}
\label{sec:program-A}
Fix one $M(h,g;u,f)$ and use $G_\ell$ as reference, without imposing
its circle symmetry on the unknown. In a chosen Bianchi or DeTurck
gauge, analyze $L_\ell=DF_{G_\ell}(0)$ on the thick product blocks,
mapping-class collars, and shear throats. Thick-region kernels come
from actual product deformations and global parameter/gauge data.
The collar must retain fibre TT transfer, constants, and small
eigenmodes; the throat must retain all mixed connection components.
Every interface uses $\Psi_e$, including holonomy.

The first target is a complete decomposition
$X_\ell=K_\ell\oplus X_{1,\ell}$ in stated mixed norms and a proved
complement inverse bound $B_\ell$. Uniform boundedness of
$\dim K_\ell$ is itself a question. The second target is the full
effective equation $\mathcal O_\ell(z)=0$, including self-interaction,
external response, and complement feedback. Only then compare
$B_\ell$, the residual, the nonlinear constants, and the positivity
radius to obtain either exact correction or a quantitative exclusion.
A rapidly growing inverse suggests moving dangerous modes into the
reference or using the receiver-weighted method; a finite invertible
matrix alone does not settle the continuous equation.

\subsection{Can cross-node or cross-edge cancellation create a true zero?}
\label{sec:program-B}
On four-node orbifolds and finite gluing graphs, preserve one parameter
space containing scale ratios, angles, finite kernels, actual background
deformations, and necessary frames. Equal local leading terms do not
imply equal global response. Quotienting by a geometric symmetry requires
proved equivariance.
If
\[
\mathcal O(t,z)=t^k\mathcal O_k(z)+\mathcal R(t,z),
\]
find actual zeros of $\mathcal O_k$ and their transversality.
A stable finite-dimensional argument needs
$\norm{\mathcal R}_{C^1}=o(t^k)$, unless same-order logarithms
are incorporated into the leading term. Nonvanishing at one fixed
$z$ excludes at most its neighborhood. A topological obstruction must
be shown to apply throughout the allowed parameter space.

A source can lie outside the receiver directions while having zero
receiver pairing; it must not be set to zero. The archived four-node
corrections distinguish ordered from physically symmetrized
polarizations (including the entries indexed $23$ and $32$).
Those labels denote tensor components, not an unexplained numerical
factor: their joint source contraction contributes a factor of two.
Their normalization and pairing rules must be restored before
using a coupling coefficient. This preserves the source correction
without treating ordered components as independent physical obstructions.

\section{Interfaces with global Einstein existence}\label{sec:interfaces}
\subsection{Fixed manifolds versus changing topology}\label{sec:changing-topology}
If $e_\sigma(M_r)=0$ for each fixed $r$, diagonal selection gives
metrics on $M_r$ with uniformly vanishing defect. It proves no
Einstein nonexistence on any $M_r$. Conversely, a topology-dependent
positive lower bound $c(M_r)\to0$ is not a construction attaining it.
For a surgery or gluing parameter, compare realized upper bounds
and valid global lower bounds in the same norm and on the same objects.
Only approximate realizability together with unrestricted nonexistence
can give a changing-topology branch counterexample.

\subsection{Exact and degenerating control families}\label{sec:controls}
Each model has useful direct checks: zero shear returns a product or
standard resolution; one node returns the actual core; fixed nonzero
shear retains its integral periods; a declared higher-order error is
smaller than its leading term. Degenerating exact Einstein families
test overstrong claims of uniform inverse bounds or vanishing local
sources. Classical Kodaira base changes give exact Einstein controls
with infinite monodromy and a normalized harmonic gap tending to zero.
Their positive signature differs from $E_D$; they invalidate an
overbroad strategy without resolving its zero-signature case.
\source{round10_bundle.md}{Kodaira controls and exact harmonic/Hodge formulas}

For computation, reconstruct Ricci and Bianchi from the full metric
rather than using the equation being checked to infer its own residual.
Before certification verify positivity, smooth endpoints, fixed
markings, and the global supremum. Interval finite-dimensional bounds
and analytic tail estimates can complement one another.

\subsection{Constructions constrain possible obstructions}\label{sec:new-obstructions}
A fixed nonzero integral shear can have vanishing Ricci cost.
An obstruction depending only on that integer and ordinary $L^p$
size must account for this fact. It may need the exact Einstein PDE,
nonlocal block response, or an unavoidable topological coupling.
Construction formulas can therefore refute an overstrong proposed
global inequality as well as produce inputs for it.

An inherited concentration selection rule makes this explicit.
For pointwise approximation with exterior convergence and vanishing
replacement volume, let $E\ge0$ be concentrated self-dual Weyl energy
and $T=2\chi+3\tau$. In the source conventions,
\[
T=T_0+\frac{E}{2\pi^2},\qquad
\alpha^2(M)\le\min\left\{\frac{V_0}{2\pi^2},
                         J_0+\frac{E}{12\pi^2}\right\},
\]
and
$T-\frac23(J_0+E/(12\pi^2))=F_0+4E/(9\pi^2)$, $F_0\ge0$.
Thus concentration alone does not move locally symmetric bulk into
the required critical monopole range. A bulk realizing the needed
geometry and saturation has not been constructed; the same
special-monopole hypotheses cannot also violate their established gap.
\source{round7_endpoint.md}{concentration, scalar cap, and the reversed slope restriction}

Likewise the inherited free-quotient construction realizes the
noncommuting full-generator bundle $E_{\rm full}$; later calculations
show that its rank-one representations factor through the base.
This group information does not become an arbitrary Einstein
obstruction.
\source{round8_multitwist.md}{fixed marked monodromy and rank-one representation calculations}

Two failed bridges remain informative. Compact convergence on a
complete model does not erase physical-period-normalized exterior
receiver flux. A growing AH core cannot be inserted into a fixed
closed topology without actual scale, volume, and boundary matching.
The construction of Section~\ref{sec:sheared} solves a compact marked
problem; it does not rely on the unproved AH-to-cusp closure.

\section{Recent external comparisons and their hypotheses}\label{sec:literature}
\subsection{Negative K\"ahler--Einstein second variations}\label{sec:Nagy}
Nagy's 2026 analysis \cite{Nagy} gives gauge-normalized second-order
deformation information on compact negative K\"ahler--Einstein
backgrounds: the expansion is determined by the squared infinitesimal
deformation and the divergence of its Kodaira--Spencer bracket.
This suggests recomputing genuine quadratic kernel responses on
product blocks. Boundary terms, noncompact domains, non-K\"ahler
backgrounds, and degenerating constants require separate analysis;
the result does not automatically remove third-order obstructions.

\subsection{Stability without a diameter bound still has geometric hypotheses}
\label{sec:Jackel}
J\"ackel's effective stability theorem \cite{Jackel} does not require
an upper diameter or volume bound, but its stated application requires
strictly negative sectional curvature, exact hyperbolic thin geometry,
Ricci derivative bounds, and both $C^0$ and $L^2$ error control.
The positive sectional curvature in \eqref{eq:ED-matrix} and the
different thin geometry of the sheared families prevent an automatic
application. One can seek a new family satisfying those conditions
or determine exactly how the current family escapes them.

\subsection{Removing singularities does not identify the input smooth structure}
\label{sec:removability}
Antonelli--Sz\'ekelyhidi's 2026 work \cite{AS} concerns removability
for $L^\infty$ Einstein metrics and RCD spaces. The four-dimensional
scope described in its abstract allows a change of smooth structure
under its codimension assumptions. A recovered Einstein geometry
advances the present problem only after its smooth type is identified
with the original input. We use this as a research comparison, not
as an unverified black-box theorem on the constructed manifolds.
The recent survey \cite{LeBrunSurvey} supplies a field map, not a
substitute for checking the assumptions and provenance of individual results.

\section{Scope of the resulting construction theory}\label{sec:scope}
Marked orbifold approximation, surface-bundle approximation with
quantitative repair obstructions, and fixed closed sheared topology
are distinct centers. Exact open models, charge identities, and
approximate connectors can support them or merit separate study.
The number of papers is not a mathematical constraint.

A stronger result should enlarge the realized smooth topologies,
admissible markings or singularities, rigorous continuous estimates,
exact correction range, or global existence information.
Every construction here carries an explicit approximation assertion
on a specified object; its unrestricted Einstein question remains
separate. On the new self-glued manifolds, topology excludes
nonnegative Einstein constants and several special negative
geometries, while completely nonsymmetric negative Einstein existence
remains unresolved; see the companion
\href{20260921_global.tex}{global manuscript}.
Resolving this negative instance alone would not defeat both
branches of \eqref{eq:branches}.

If a route is excluded, retain its valid formulas and failure mechanism
and permit a different topology, background, all-mode ansatz, or method.
If an exact metric is found, that is a positive geometric result.
Finite-order cancellation must still be distinguished from actual
integrability of the full equation.

\appendix
\section{Proof navigation and inherited boundaries}\label{app:sources}
This article integrates the dated construction synthesis with the
archived twelve-round proof material. Baseline manuscripts remain in
this directory; original round files remain under
\nolinkurl{../references/rounds/}. They retain full parameter
dependencies, reviews, scripts, and historical intermediate arguments.
This article compiles without reading those files, but it does not
replace their complete proofs. No fresh execution of every historical
program or comprehensive priority determination is reported.

\small
\begin{longtable}{L{.34\textwidth}L{.59\textwidth}}
\toprule Source & Content and current use\\\midrule\endhead
\href{20260916_construction.tex}{\nolinkurl{20260916_construction.tex}}
& Marked $A_1$ core/annulus/exterior proof; balanced $E_D$; complete-kernel and four-node calculations. Inherited constructions, not new-round claims.\\
\href{20260912.tex}{\nolinkurl{20260912.tex}}
& Flux and ALE extensions; fixed $E_D$, all Ricci blocks, geometric bounds, three-scalar correction, global pairings.\\
\href{20260916_global.tex}{\nolinkurl{20260916_global.tex}}
& Geometric free-quotient monodromy and already-existing charged Koiso germs.\\
\href{../references/rounds/construction.md}{\nolinkurl{construction.md}}
& Exact polarized base-TT cancellation and correction scales; later whole-fibre arguments remove the dominant-mode restriction.\\
\href{../references/rounds/independent_review.md}{\nolinkurl{independent_review.md}}
& Contractions, geometric Taylor remainders, background measure, and distinction between cubic order and an optimal coefficient.\\
\href{../references/rounds/round2_construction.md}{\nolinkurl{round2_construction.md}}
& Ten blocks, bounded nonseparable kernel, invariant $L^2$ gap, exact cylinder parametrization, quadratic sources, physical scales, and the collapsing-normalization failed shortcut.\\
\href{../references/rounds/round2_construction_review.md}{\nolinkurl{round2_construction_review.md}}
& Kernel scope, linear correction, normalization, and open versus closed objects.\\
\href{../references/rounds/round3_construction.md}{\nolinkurl{round3_construction.md}}
& Whole nonlinear source, weighted inverse, moving gauge, contraction, Bianchi removal, end models, second-variation flux, and prior-work boundaries.\\
\href{../references/rounds/round3_construction_review.md}{\nolinkurl{round3_construction_review.md}}
& Weighted coercivity, complete target space, gauge removal, and end decay beyond formal Ricci expansion.\\
\href{../references/rounds/round3_matching_next.md}{\nolinkurl{round3_matching_next.md}}
& End shifts, conditional matching, exterior adjoint defect, physical periods, and the actual horizontal closed detector; no invented vertical detector.\\
\href{../references/rounds/round4_construction.md}{\nolinkurl{round4_construction.md}}
& Nonlinear charge, positive finite response matrix, constant-twist gauge, boundary energy, and higher coefficients of the same mechanism.\\
\href{../references/rounds/round4_exact_core.md}{\nolinkurl{round4_exact_core.md}}
& Unchanged exact-core insertion obstruction through cocompactness and analytic continuation.\\
\href{../references/rounds/round4_review_construction.md}{\nolinkurl{round4_review_construction.md}}
& Fixed-radius decay, arbitrary cocompact groups, and no quantitative small-repair claim.\\
\href{../references/rounds/round5_construction.md}{\nolinkurl{round5_construction.md}}
& Whole-fibre tightness, arbitrary perturbations, weak quadratic convergence, bounded gauge, and the superseded $\varepsilon^{3/2}$ exclusion radius.\\
\href{../references/rounds/round5_receiver_independent.md}{\nolinkurl{round5_receiver_independent.md}}
& Independent complete-kernel and determinant-preserving signed receiver.\\
\href{../references/rounds/round5_construction_rigidity_review.md}{\nolinkurl{round5_construction_rigidity_review.md}}
& Tight detector mass, justified $C^1$ convergence, gauge response, and exterior terms; no unjustified strong $C^2$ convergence.\\
\href{../references/rounds/round5_rank_obstruction.md}{\nolinkurl{round5_rank_obstruction.md}}
& Rank-one receiver obstruction in its stated repair class, not unrestricted Einstein nonexistence.\\
\href{../references/rounds/round5_modulation.md}{\nolinkurl{round5_modulation.md}}
& Logarithmic approximate TT mode defeating maximum-norm modulation; later replaced by the sufficient receiver-weighted inverse.\\
\href{../references/rounds/round6_weighted_inverse.md}{\nolinkurl{round6_weighted_inverse.md}}
& Nonseparating scalar gap, all weighted blocks, partner factorization, actual closed-fibre transfer, gauge, and signed quadratic response.\\
\href{../references/rounds/round6_construction.md}{\nolinkurl{round6_construction.md}}
& Linear-size closed exclusion, finite collars, angular error, nonzero gauge forcing, all same-order errors, and limits beyond the small fixed multiple.\\
\href{../references/rounds/round6_receiver_review.md}{\nolinkurl{round6_receiver_review.md}}
& Weighted domains and the term $\mathcal L_XF$ that cannot be discarded.\\
\href{../references/rounds/round6_global.md}{\nolinkurl{round6_global.md}}
& Refinement of charged germs; local support-volume limitation; one-ended $D^2\times T^2$ charge release; conditional end flattening, not closed gluing.\\
\href{../references/rounds/round7_endpoint.md}{\nolinkurl{round7_endpoint.md}}
& Concentration identity, scalar cap, critical ratio, reversed KE slope, and missing saturated bulk.\\
\href{../references/rounds/round8_multitwist.md}{\nolinkurl{round8_multitwist.md}}
& Explicit fixed free-quotient bundles with inherited interpolation, rank-one curve killing, representation volumes, and full-generator examples.\\
\href{../references/rounds/round9_connector.md}{\nolinkurl{round9_connector.md}}
& Exact free-circle smooth-boundary obstruction; not a theorem for double-cylinder product asymptotics.\\
\href{../references/rounds/round10_kottler.md}{\nolinkurl{round10_kottler.md}}
& Smooth fixed double, exact Ricci eigenvalues, exact exteriors, optical localization, failed symmetry-extension hypothesis, and approximate non-gauge mode.\\
\href{../references/rounds/round10_cusp_doubling.md}{\nolinkurl{round10_cusp_doubling.md}}
& Li--Liu source hypotheses, global bundle connection, complete central Ricci blocks, relative two-jets, seam error, Chern flux, and cusp-type restrictions.\\
\href{../references/rounds/round10_bundle.md}{\nolinkurl{round10_bundle.md}}
& Classical Kodaira base changes and harmonic/Hodge formulas; exact control family with a different signature from $E_D$.\\
\href{../references/rounds/round11_closed_gluing.md}{\nolinkurl{round11_closed_gluing.md}}
& Finite marked blocks, real collars, path and shear Ricci blocks, all-jet assembly, exact volume, almost parallel curvature, and both-orientation almost K\"ahler structures.\\
\href{../references/rounds/round11_graph_topology.md}{\nolinkurl{round11_graph_topology.md}}
& HNN group, nonextension of integral shear, examples beyond all oriented surface bundles, and inapplicable infranil graph obstructions.\\
\href{../references/rounds/round12_repair.md}{\nolinkurl{round12_repair.md}}
& Exact scalar constraint, positive receiver, marked displacement, transport backreaction, spatial-radius escape, and common-area rigidity.\\
\href{../references/rounds/round12_symmetry.md}{\nolinkurl{round12_symmetry.md}}
& Discrete universal-cover isometries and total local symmetry breaking; no quantitative asymmetric gap.\\
\href{../references/rounds/round6_verification.md}{\nolinkurl{round6_verification.md}},
\href{../references/rounds/round10_verification.md}{\nolinkurl{round10_verification.md}}
& Symbolic, coordinate-Ricci, scaling, Kottler, and cusp checks; no continuous PDE certificate.\\
\href{../references/rounds/round11_verification.md}{\nolinkurl{round11_verification.md}},
\href{../references/rounds/round12_verification.md}{\nolinkurl{round12_verification.md}}
& Full path/shear blocks, deliberately wrong formulas rejected, topology consistency, and the spatial-lapse scalar identity within its exact scope.\\
\bottomrule
\end{longtable}
\normalsize
The original \nolinkurl{*check*.py} files and recorded reproduction
outputs remain associated with their proof versions. Finite algebra,
integer topology bookkeeping, sampling, continuous spectral tails,
nonlinear remainders, and global PDE arguments have different roles.
The archived additional checks of OT topology, specific chirality
applications, negative GH separation, and independent coordinate
calculations remain distinct from failed problem-W splicing and a
specified cutoff's $1/3$ leading term. Historical reviews and source
files are preserved, not overwritten by this article.

\begin{thebibliography}{99}
\bibitem{OzuchII} T.~Ozuch, \emph{Noncollapsed degeneration of Einstein
4-manifolds II}, Geometry \& Topology (2022).
\url{https://arxiv.org/abs/1909.12960}. See the pointwise Ricci
approximation discussion and logarithmic cutoffs.
\bibitem{OzuchInt} T.~Ozuch, \emph{Integrability of Einstein deformations
and desingularizations}. \url{https://arxiv.org/abs/2105.13193}.
\bibitem{Biquard} O.~Biquard, \emph{D\'esingularisation de m\'etriques
d'Einstein. I}, Inventiones Mathematicae (2013).
\url{https://arxiv.org/abs/1105.4904}.
\bibitem{OR} P.~Orlik and F.~Raymond, \emph{Actions of the torus on
4-manifolds. I}.
\url{https://webhomes.maths.ed.ac.uk/~v1ranick/papers/orlikraymond1.pdf}.
\bibitem{Rupflin} M.~Rupflin and P.~Topping, \emph{Horizontal curves
of hyperbolic metrics}, Appendix A.
\url{https://arxiv.org/abs/1605.06691}.
\bibitem{Ji} L.~Ji and R.~Wentworth, \emph{Spectral convergence on
degenerating surfaces}. \url{https://math.umd.edu/~raw/papers/spectrum.pdf}.
\bibitem{Nagy} P.-A.~Nagy, \emph{Einstein deformations of K\"ahler
Einstein metrics}, 9 March 2026.
\url{https://arxiv.org/abs/2603.09028v1}.
\bibitem{Yao} Y.~Yao, \emph{Resolution of compact Einstein orbifolds
in general dimensions}, version 2, 5 July 2026.
\url{https://arxiv.org/abs/2603.14929v2}.
\bibitem{HO} Y.~Huang and T.~Ozuch, \emph{Regularity of Einstein
5-manifolds via 4-dimensional gap theorems}.
\url{https://arxiv.org/abs/2512.21317}.
\bibitem{Jackel} F.~J\"ackel, \emph{Effective stability of negatively
curved Einstein metrics in dimensions at least 4}.
\url{https://arxiv.org/abs/2502.16692}. The comparison retains the
sectional-curvature, thin-region, derivative, and error hypotheses.
\bibitem{AS} G.~Antonelli and G.~Sz\'ekelyhidi, \emph{Removability of
non-isolated singularities for Einstein metrics and RCD spaces},
September 2026. \url{https://arxiv.org/abs/2609.01464};
\url{https://cvgmt.sns.it/paper/7959/}.
Used only as the qualified comparison in Section~\ref{sec:removability}.
\bibitem{LiLiu} M.~Li and H.~Liu, \emph{Poincar\'e--Einstein
4-manifolds with cusps}, 25 May 2026.
\url{https://arxiv.org/abs/2605.25462v1}.
\bibitem{FW} B.~Farb and S.~Weinberger, \emph{Isometries, rigidity
and universal covers}, Annals of Mathematics 168 (2008).
\url{https://math.uchicago.edu/~farb/papers/isoms.pdf}.
\bibitem{Wang} Q.~S.~Wang, \emph{Computer-assisted construction of
$SU(2)$-invariant negative Einstein metrics}, Annals of Global Analysis
and Geometry 69 (2026), no.~3, Paper 18.
\url{https://arxiv.org/abs/2504.21644v2}.
\bibitem{LeBrunSurvey} C.~LeBrun, \emph{An Overview of Einstein
4-Manifolds}, 2 September 2026.
\url{https://arxiv.org/abs/2609.03033}.
\end{thebibliography}
\end{document}
